University of Wollongong
Research Online
Faculty of Informatics - Papers (Archive) Faculty of Engineering and Information
Sciences
1-1-2007
Relating business process models to goal-oriented requirements models in
KAOS
George Koliadis
University of Wollongong, gk56@uowmail.edu.au
Aditya K. Ghose
University of Wollongong, aditya@uow.edu.au
Follow this and additional works at: https://ro.uow.edu.au/infopapers
Part of the Physical Sciences and Mathematics Commons
Recommended Citation
Koliadis, George and Ghose, Aditya K.: Relating business process models to goal-oriented requirements
models in KAOS 2007.
https://ro.uow.edu.au/infopapers/573
Research Online is the open access institutional repository for the University of Wollongong. For further information
contact the UOW Library: research-pubs@uow.edu.au 
Relating business pr Relating business process models t ocess models to goal-oriented r o goal-oriented requirements models in KA ements models in KAOS
Abstract
Business Process Management (BPM) has many anticipated benefits including accelerated process
improvement, at the operational level, with the use of highly configurable and adaptive “process aware”
information systems [1] [2]. The facility for improved agility fosters the need for continual measurement
and control of business processes to assess and manage their effective evolution, in-line with
organizational objectives. This paper proposes the GoalBPM methodology for relating business process
models (modeled using BPMN) to high-level stakeholder goals (modeled using KAOS). We propose
informal (manual) techniques (with likely future formalism) for establishing and verifying this relationship,
even in dynamic environments where essential alterations to organizational goals and/or process
constantly emerge.
Disciplines
Physical Sciences and Mathematics
Publication Details
This conference paper was originally published as Koliadis, G and Ghose, AK, Relating Business Process
Models to Goal-Oriented Requirements Models in KAOS, in Proceedings of the 2006 Pacific-Rim
Knowledge Acquisition Workshop (PKAW-2006), Lecture Notes in Computing Science, 2007.
This conference paper is available at Research Online: https://ro.uow.edu.au/infopapers/573 
Relating Business Process Models to
Goal-Oriented Requirements Models in KAOS
(Preprint)
George Koliadis1 and Aditya Ghose1
School of Information Technology and Computer Science
University of Wollongong
Wollongong, N.S.W. Australia
{gk56, aditya}@uow.edu.au
http://www.uow.edu.au
Abstract. Business Process Management (BPM) has many anticipated
benefits including accelerated process improvement, at the operational
level, with the use of highly configurable and adaptive “process aware”
information systems [1] [2]. The facility for improved agility fosters the
need for continual measurement and control of business processes to
assess and manage their effective evolution, in-line with organizational
objectives. This paper proposes the GoalBPM methodology for relating
business process models (modeled using BPMN) to high-level stakeholder
goals (modeled using KAOS). We propose informal (manual) techniques
(with likely future formalism) for establishing and verifying this relationship, even in dynamic environments where essential alterations to
organizational goals and/or process constantly emerge.
1 Introduction
Business Process Management (BPM) in its “third wave” [1] has been conveyed
as: enabling intelligent business management [3]; facilitating the redesign and
organic growth of information systems [4]; and, obliterating the business - IT
divide [1]. Business processes undergo an evolutionary life-cycle of change. This
change is brought on by the need to satisfy the constantly changing goals of
varied stakeholders and adapt to the accelerating nature of change in today’s
business environment [5]. The need for change is best described in [6] as the
transition from an initial “unsatisfactory” (i.e. as-is) state to a new hypothetically “desired” (i.e. to-be) state. The desired state is theoretically based on the
assumption that it more effectively satisfies related operational goals [4] [7] [8]
in-line with higher-level strategic goals. It is therefore important that the criterion for effective process change - i.e. stakeholder goals, be explicitly stated,
communicated and traceable to any changes that are proposed, approved, and/or
implemented.
The new-found agility provided by BPM, however presents the need for methods to successfully control and trace the evolution of processes. This need is affirmed in [7], by stating that organizations evolve from their original intentions
2 George Koliadis et al.
through complex and unpredictable growth. BPM aims to support the evolution
of organizations and their processes, however controls are still needed to ensure
that operational as well as higher-level goals (i.e. of more strategic concern) are
continually satisfied, allowing for “organizational growth in the right direction”.
In order to meet goals however, there is a need to support traceability between
processes and organizational goals - “You can’t manage what you can’t trace”
[9].
We have proposed a method (GoalBPM) to support the controlled evolution of business processes. Control is supported through the explicit modeling of
stakeholder goals, their relationships (be it either refinement, conflict or obstruction), and their evolution traceable to related business processes. GoalBPM is
used to couple an existing and well-developed, informal-formal goal modeling and
reasoning methodology - i.e. KAOS [10], and a newly developed business process
modeling notation - i.e. BPMN [11]. This is achieved through the identification
of a satisfaction relationship between the concepts represented. GoalBPM itself
can be seen as an “adapter” that integrates the two models, to support their
co-evolution and synergistic use.
This paper firstly presents a background to the associated domains of business process modeling and goal-oriented requirements engineering. An informal
overview of the GoalBPM method is subsequently outlined with a simple example for illustration.
2 Business Process Modeling with BPMN
We have initially chosen the Business Process Modeling Notation (BPMN), developed by the Business Process Management Initiative (BPMI.org) for use in
the construction of GoalBPM.
A Business Process is a set of dynamically co-ordinated activities controlled
by a number of dependent, social participants. Processes are represented in
BPMN using flow objects: events (circles), activities (rounded boxes), and
decisions (diamonds); connecting objects: control flow links (unbroken directed lines), and message flow links (broken directed lines); and swim lanes:
pools (high level boxes containing a single process), and lanes within pools (subboxes).
We refer to Figure 2, a public Package Sorting process, as an example to
illustrate BPMN. The process requires the interaction of two high-level process
participants - the Transport Organization, and a Transport Authority. Collaboration between participants on the model is represented by message flow links
between activities within pools. Responsibility within the Transport Organisation is delegated to two roles - the Sort Operations, and a Bond Operations.
Responsibility assignment within a pool is represented using lanes (i.e. pool divisions). Each pool within a process model represents a single process. Processes
are initiated by a start event, represented as a circle at the beginning of each
pool. Control flow links between activities, decisions and events, represent the
controlled progression through each process. A decision gateway, (i.e. diamond)
Supporting BPM with GORE 3
can be seen in the figure, identifying the need to make a choice on whether to
bond a package. Finally, the process is completed with an end event, or bold
circle toward the end of a process.
3 Goal-Oriented Requirements Engineering with KAOS
We have chosen a GORE method that is focused toward both the early and late
phases of RE, specifically KAOS (Knowledge Acquisition in autOmated Specification of software systems) [10], to represent the organizational goals related to
business process execution.
Goal Declaration in KAOS Goals are declared in terms of desired, timely
effects within a composite system (e.g. Achieve[MeetingScheduled]). Goals are
conceptually modeled in KAOS on a semantic net that represents a hierarchy
of parent goals and their refinements into sub-goals. Goals that exist higher
in the hierarchy represent the high-level goals (i.e. strategic concerns) of the
organization. These goals are not “clear cut”, in the sense that their satisfaction
is complex and cannot be proven without common interpretation. These goals
are then refined down the hierarchy into sub-goals that are more operational
in nature. That is, their assignment to a small group of individuals responsible
for a number of operational activities illustrates the means by which they are
satisfied.
Goals can be either ‘AND’ or ‘OR’ refined (Figure 3). An ‘AND’ refinement
of a goal states that the parent goal is satisfied if all the goals in the refinement
are satisfied. An ‘OR’ refinement on the other hand states that the parent goal
is satisfied if a single refinement is satisfied. This allows for the modeling of alternative refinements for goal satisfaction. KAOS also provides a criterion for
halting the refinement process, in that if a goal can be assigned the sole responsibility of a single environmental role (i.e. agent in the composite system); there
is no need for further goal refinement to occur. This also provides a means by
which to make the transition between goals and the constrained operations that
satisfy those goals [12].
Goal Definition in KAOS KAOS supplies an optional formal assertion layer
that allows for the specification of goals in Real-Time Linear Temporal Logic
(RT-LTL) [10] [13]. These formal goal assertions allow for precise specification of
goals, as well as supporting the use of developed formal reasoning techniques that
aid in identifying/resolving conflicts between goals and proving absolute/partial
goal satisfaction.
A formal goal definition in KAOS begins with the assertion of the objects
the goal concerns. In KAOS, these objects are declared in the object model.
The definition then states the desired temporal ordering of states the concerned
objects must hold in order to satisfy the goal.
Goals are defined in the form:
4 George Koliadis et al.
Fig. 1. Modeling goals in KAOS
C ⇒ opT,
where C and T are assertions about environmental situations (i.e. current
and target), and op is a temporal operator that signifies the desired temporal
nature of the target situation (T), in relation to a current situation (C). This
paper requires knowledge only of the two most used operators. Namely, ‘at some
time in the future’ denoted by a open diamond ‘’, and ‘at all times in the
future’ denoted by an open square ‘’. For a complete list of the operators used
in KAOS, see [10].
Take for example, the ‘PackageSortedToDestination’ goal below:
Goal Achieve[PackageSortedToDestination]
InformalDef If a package is received at a sort facility, then the package
will eventually be forwarded to its known destination.
FormalDef ∀p: Package, sf: SortFacility
Received(p, sf) ⇒ Forwarded(p, p.Destination)
Also note: package destinations (i.e. p.Destination) are attributes of a package.
Patterns for Declaring and Defining Goals KAOS defines a number of commonly used “Goal Patterns” that generalize the timeliness of target situations.
They provide an informal method to initially declare goals, as well as to guide formal definition. Achieve Goals (C ⇒ T) desire achievement ‘some time in the future’. That is, the target must eventually occur. (e.g. Achieve[PaymentRecieved]).
Cease Goals (C ⇒  − T) disallow achievement ‘some time in the future’. That
is, there must be a state in the future where the target does not occur (e.g.
Cease[Operation]). Maintain Goals (C ⇒ T) must hold ‘at all times in the future’ (e.g. Maintain[EmployeeSafety]). Avoid Goals (C ⇒ −T) must not hold
‘at all times in the future’ (e.g. Avoid[LateEntry]).
Supporting BPM with GORE 5
4 Linking Goal and Process Models
The GoalBPM methodology relies on establishing relationships between goal
models and process models. This relationship is established in two stages. First,
traceability links are established between goal nodes in the goal model to activities (or end events of complete processes) in the process model. Second,
satisfaction links are established between goals and processes.
A traceability link is an informal statement of a relationship between a goal
and a process (or sub-process). In establishing such a link, we are effectively
asserting that the goal in question has some bearing on the process (or subprocess) under consideration. A traceability link does not necessarily lead to a
satisfaction link. Sometimes, a process or sub-process may be related to a goal
because it obstructs it, in the sense of [14]. We draw a traceability link between a
goal and a process end event if the entire process in question has some bearing on
the goal. We draw a traceability link between a goal and an activity if the subprocess ending in that activity relates to the goal. In general, traceability links
need to be established by analysts. Some guidance can be offered in this process
by using cues present in the goal and process models. For instance, the names
of goals and processes (or activities) can suggest traceability links. A transport
organization’s operational goal to achieve “PackageSortedToDestination” can be
traceable to the “Package Sorting” process within the organization.
Traceability between goals and processes can be identified through cross examination of the links between the pre/post conditions for specific processes
and the pre/post conditions for specific goals. A process is made available for
execution when a specific pre condition has been met (e.g. a customer that has
submit a registration form in the ‘Register New Customer’ process), and completed upon meeting a post condition (e.g. the customer is validated and their
details are stored). These pre and post conditions can be related to the pre and
post conditions for specific organizational goals. Take for example the goal, ‘all
new customer registrations require credit reporting and verification’, is related
to the ‘Register New Customer’ process by way of the pre condition.
A satisfaction link is a traceability link where the process (or sub-process) in
question satisfies the goal involved in the link. A satisfaction link can be of two
types:
– Normative satisfaction links: These indicate that a process or sub-process
must satisfy the relevant goal. Such links articulate desired states of affairs.
– Descriptive satisfaction links: These indicate the “as-is”. We obtain descriptive satisfaction links from effect annotations of processes, using techniques
that we discuss in the next section.
5 Using Model Annotations to Verify Goal Satisfaction
GoalBPM establishes satisfaction links between goals and process models in
three steps. First, it annotates process models with effect annotations. Second, it
6 George Koliadis et al.
identifies a set of critical trajectories from a process model. Third, it identifies the
subset of the set of traceability links that represent satisfaction links by analyzing
critical trajectories relative to process effect annotations. The satisfaction links
thus obtained are descriptive satisfaction links. A final step in GoalBPM is to use
a comparison of the set of normative satisfaction links with the set of descriptive
satisfaction links to drive the processes of goal model update and/or process model
update.
Our approach may be viewed as an instance of the state-oriented view [15]
[16] [17] of business processes as opposed to the agent-oriented or workflow views.
However, we are not explicitly state-based in that we do not seek to obtain state
machine models from process models, for two reasons. First, BPMN models
in general do not guarantee finite state systems, making the application model
checking techniques difficult. Second, the derivation of state models from BPMN
models appears difficult at this time, due to the high-level, abstract nature of
BPMN models.
5.1 Effect Annotations
A process activity (i.e. as represented in BPMN as a rounded box) is an element
on a process model that indicates required state transitions in order for the
process to progress toward the achievement of all related goals. The labeling
of an activity (e.g. ‘Register New Customer’) generalizes a number of possibly
desired/undesired results, or outcomes. Due to this generalization, most process
models do not satisfactorily depict the lowest level state achievements required
for process progression. They are too ‘high-level’ and do not provide an in depth
understanding of the process and its ability to achieve desired goals or objectives.
This understanding is important when trying to prove goal achievement, which
is reliant on the achievement of certain target states. In order to provide greater
understanding of process models (i.e. state transitions in particular) to support
their analysis in relation to the goals they hope to satisfy, we augment the process
model with ‘effect annotations’.
An effect is the result (i.e. product or outcome) of an activity being executed
by some cause or agent. It indicates the achievement of a certain environmental
state communicated through an event. An effect annotation relates a specific
result or outcome to an activity on a business process model. It explicitly states a
result of the activity if the conceptual model were to be hypothetically executed.
A cause relationship exists between a process activity and an effect (i.e. the
process activity causes the effect to occur). An activity can cause many effects
and an effect can be caused by a number of activities.
The manner in which an activity is executed may result in alternative outcomes. For this reason, effect annotations are related to activities in an AND/OR
refinement similar to goal refinement in KAOS. This alternative execution is derived from the dynamic co-ordination represented in the process model through
decision gateways (i.e. diamonds in BPMN). Decisions commit to alternate paths
of execution based on the current state of the process. This is achieved through
the specific effects of prior activities that have been executed. The conditions for
Supporting BPM with GORE 7
the decision on which choice of path to commit to can help to identify important
effects on prior activities in the current, or in other processes. These influential
activities and their required effects for the current path of execution, need to be
identified and represented along with the effects of the current process to prove
goal satisfaction.
We define an effect annotation to include:
– a label that generalizes the behavior of the effect in relation to its environment (e.g. ‘CustomerDetailsStored’). Whereas the labeling of an activity is
made in the optative mood (i.e. a desire), an effect annotation is made in
the indicative mood (i.e. a fact).ling of an activity is made in the optative
mood (i.e. a desire), an effect annotation is made in the indicative mood (i.e.
a fact).
– a designation specifying whether the effect is a ‘normal’ (i.e. expected) outcome for the activity that in turn aims toward goal achievement, or an ‘exceptional’ (i.e. unaccepted) effect that deviates from goal achievement. (e.g.
‘RegistrationValidated’ may be a normal outcome for a customer registration
activity, whereas ‘RegistrationRejected’ may be exceptional)
– a informal definition an informal definition describing the effect in relation
to the result achieved in its environment (e.g. ‘The details relating to the
current customer have been stored within the system.’). This provides an
informal explanation (i.e. meaning) of the effect in relation to the real-world
environment.
– a formal definition (optional) defining achieved states to aid in mapping to
formal goal definitions in the chosen goal definition formalism (i.e. in this
case KAOS). (e.g. ‘∀ c: Customer, (∃ cr: CustomerRecord) Stored(c.Details,
cr)’)
At the tool level, effect annotations can be viewed on a business process model
graphically, or added to meta-information relating to the process activities. They
can then be analyzed along with the process and associated goals as described
in the subsequent sections.
5.2 Trajectory Decomposition
The dynamic co-ordination of activities controlled by responsible agents is represented / supported on a business process model by way of decision gateways (i.e.
diamonds in BPMN). This manner of ‘per instance’ process control allows for
the existence of many process trajectories through the process (i.e. possibly even
an infinite number when cycles are included). We classify a single trajectory as
a unique and supported sequence of activity execution. Each trajectory results
in a ‘cumulative effect’ for the given process. We use the term trajectory to signify a specific ‘chosen path’ through the process that results in a specific/unique
outcome or ‘cumulative effect’.
During the trajectory decomposition process, specific effects need to be chosen where alternative effects are available on each activity. This choice relies
8 George Koliadis et al.
on the decisions influencing the path of the particular trajectory. We choose
between alternative effects based on their conformance to the current trajectory.
Business process models support and represent exceptional trajectories. These
trajectories do not necessarily satisfy process goals. They react to exceptional
events that occur in the process by re-routing the current path of execution so
that alternative steps can be taken to either resolve the exceptional situation or
abort the process. We can therefore label a trajectory as having either a normal
or exceptional type, based on its final achieved state. Commonly, a trajectory
that cannot resolve exceptions and requires termination prior to meeting all the
goals it must satisfy, is classed as exceptional. This is discussed in the following
section.
The identification of all unique process trajectories can be difficult, given
the complex nature of activity interleaving and iteration possible in the process
models. For this purpose, it is recommended that an automated method for
deriving all possible process trajectories be available.
5.3 Goal Satisfaction and Cumulative Effect Assessment
We firstly progress through each process trajectory and compare effects with
traceable satisfaction goals in the goal model. Effects are compared to the desireability and temporal ordering of effects in normative goals and descriptive
satisfaction links are established as we progress through the trajectory.
We then analyze and classify each trajectory as either normal or exceptional.
A normal trajectory in relation to the goal model leads to the satisfaction of all
normative goals. An exceptional trajectory, as described in the previous section,
satisfies a limited number of normative satisfaction goals.
In order for a satisfaction relationship to exist between a goal model and an
associated process model, there must be at least one normal trajectory. That is,
the process model must support at least one valid means by which to satisfy the
required normative goals of the process.
Finally, we analyze the outcome of the satisfaction process, identifying whether
the process supports the achievement normative satisfaction goals and classify
the satisfaction relationship between the process and the associated goal model
as either strong, weak, or unsatisfied. A strong satisfaction relationship is determined if all possible trajectories are ‘normal’ (i.e. satisfy all associated goals).
On the other hand, a weak satisfaction relationship is said to exist when there is
at least one ‘exceptional’ trajectory and one ‘normal’ trajectory. This classification, delineating between weak and strong satisfaction, can be important when
evaluating the competency of the process in recovering from exceptional situations that may arise during enactment. An unsatisfied satisfaction relationship
is the result of there not being a single ‘normal’ trajectory decomposed from the
process. This classification requires that changes are made to either the process
and/or goal model to establish a weak or strong satisfaction relationship.
Supporting BPM with GORE 9
6 Example
We apply the GoalBPM to a single case within a Transport Organization for
illustration. GoalBPM is specifically applied to a core operational ‘Package Sorting’ process within the organization (see Figure 2). Furthermore, we introduce
change to the goal model, and identify inconsistencies with the current business
process that need to be addressed to maintain the satisfaction relationship.
6.1 Current Business Context and Process
Fig. 2. The ’Package Sorting’ Process of the Transport Organization
The organisation has recently chosen to undergo redesign to maintain effectiveness specifically in relation to increasing regulatory requirements including
the screening of packages. This has impacted the internal package sorting operations of the organisation by introducing the requirement for every package to be
screened by a representative of the authority upon arrival and prior to it leaving
the sort facility for delivery or routing.
The operational objective the ‘Package Sorting’ process aims to achieve is
the prompt routing of packages, upon arrival to a sorting facility, to their respective destinations. There are three primary participants in the process whose
objectives must be met: Transport Organization - whose concern is the efficient
routing of packages to their destinations by assigning responsibility to internal
sort operations as well as bond operations whose role is to liaise with transport
authorities for prompt delivery clearance; Customers - whose concern is prompt
delivery and package traceability; and Transport Authorities - whose concern is
with maintaining a high level of integrity in regards to border control through
package screening. These requirements are represented on the goal model in
Figure 3, with definitions supplied in Figure 4.
10 George Koliadis et al.
The ‘Package Sorting’ process in Figure 2, represents the current ‘as-is’ coordination of activities and interactions aimed toward achieving traceable goals.
We apply GoalBPM to the process and goal model in order to evaluate the current state of the satisfaction relationship before introducing the aforementioned
changes to the goal model.
Fig. 3. Goals Traceable to the ‘Package Sorting’ Process
6.2 Applying GoalBPM
We apply the proposed GoalBPM to prove the satisfaction relationship between
the goal model and the process model. We specify effects as they would be
defined formally due to space limitations (e.g. Arrives(p, sf ) is equivalent to
PackageArrivesAtTheSortFacility, as can be inferred from the goal definitions in
Figure 4).
Analyzing Traceability We firstly declare boundaries to help guide an evaluation of the traceability relation by identifying the required pre and possible post conditions for the ‘Package Sorting’ process. Pre conditions include:
Arrives(p, sf) AND Provided(p.DeliveryDetails, ta). Post conditions include:
Sorted(p, p.Destination) OR Held(p, ta).
The initial pass at goal traceability for the ‘Package Sorting’ process identified
three high-level goals. Further analysis identifies the specific refinements that are
required for satisfaction at some point during process enactment. These goals
are declared in Figure 3 and defined in Figure 4.
Supporting BPM with GORE 11
Fig. 4. Definitions for Traceable ‘Package Sorting’ Goals
Fig. 5. Tabulated Effect Annotation, Trajectory Decomposition, and Goal Satisfaction
12 George Koliadis et al.
Effect Annotation Firstly, we annotate the model with effects to identify the
achievable (and alternative) outcomes of activities in the current process. We also
include the pre-conditions themselves, and any other relevant/influential effects
that may be caused by other processes that have a direct impact on process
decisions and coordination. These annotations are listed in Figure 5.
Process initiation is governed by two conditions: the arrival of packages to the
sort facility and the provision of package information to transport authorities.
It is also identified that the prior provision of information to authorities may
allow for the rapid clearance of packages for delivery prior to the sorting process
initiating. This may occur if the requirements of the authority can be identified
as being met. These effects are also added to the list of relevant/influential effects
that may have occurred prior to process initiation.
Each activity is then analyzed and annotated with normal and exceptional effects. Scanning a package results in the delivery details being known. The package
may be bonded for clearance with another scan being applied, and alternatively
held if the transport authority requests. The latter is an exceptional effect that
occurs due to the package meeting some characteristics that require it to be
given to the authorities that then take sole control of the package from that
point on. The outcome of releasing a package is the passing of transport authority requirements and ultimately clearance. The sorting activity then results in
the package being sorted to its destination. The final scanning activity results in
another update of package location.
Trajectory Decomposition We now decompose trajectories. An analysis of
the conditions influencing the choice of paths at decision gateways also guides
the choice of prior and/or alternative effects when accumulating them for any
given trajectory. We identify that that there are three trajectories represented
in the process model. These are listed in Figure 5 with the actions that compose
the trajectory, and the effects chosen from associated annotations.
The first trajectory decomposed from the model represents the prior clearance of a package, resulting in the eventual sorting of the package to its destination. The second trajectory results in packages requiring bonding prior to
clearance, however the exceptional alternative is selected based on the adjacent
decision gateway, and the package is held by the transport authority. The final
trajectory is categorized by requiring the package to be bonded however, the package is eventually cleared by passing the authorities requirements, as well as being
sorted to, and scanned at its destination.
Satisfaction Analysis We iterate through each trajectory, accumulate effects,
and correlate accumulations with the desired effects and their temporal ordering
on the goal model. Goals must be satisfied by firstly achievement of their pre
conditions (i.e. the antecedent) with the subsequent achievement of the consequent in conformance to the temporal pattern chosen. Goal definitions are listed
in Figure 4, with the results of the satisfaction process listed in Figure 5.
Supporting BPM with GORE 13
The satisfaction process identifies two normal trajectories and one exceptional trajectory representing a weak satisfaction relationship between the goal
model and the process model. The achievement of the operational objective to
sort packages to their destinations occurs in trajectories (1) and (3), however it
is not achieved in (2) due to an exceptional result occurring when the package
is bonded (i.e. it is held).
6.3 Changes to the Goal Model
We now introduce the newly acquired regulatory requirements for package screening that were mentioned previously. These requirements are added to the required goals for satisfaction by the process.
Goal Alterations Required alterations to the goal model, whether they be
the addition, removal or modification of goals, ultimately result in modification
to the desirability and/or temporal ordering of effects. Take for example, the
newly acquired requirement for package screening, as represented on the adjusted
portion of the goal model in Figure 6. This has introduced a newly desired effect
requiring all packages be screened by the transport authority once they arrive at
a sort facility, or formally:
∀p: Package, sf: SortFacility, ta: TransportAuthority
Arrives(p, sf) ⇒ Screened(p, ta)
Additions of new effects to goal models frequently impact other goals and
subsequently their refinements. As noted on Figure 6, the desired Screened(p, ta)
is also added to the pre conditions for PackageHeld, as well as the PackageCleared
goal and its refinements.
Process Implications and Evolution Any alterations to the goal model will
proportionately affect the desired achievement or coordination of effects within
the business processes that are assigned their operationalization. We re-evaluate
the satisfaction relationship between the ‘Package Sorting’ process and its goals,
and apply some informal analysis to identify specific changes required at the
process level.
Upon evaluation of the satisfaction relationship, it is identified that previously normal trajectories (i.e. (1) and (3)), are now also exceptional due to their
inability to satisfy regulatory requirements. This is consistent with the modifications to the goal model.
Upon further analysis of unsatisfied goals, we deduce that the Screened(p,
ta) effect is required at some time after package arrival and prior to package
clearance including bonding, as well as package holding. A decision is made as
to the addition of the desired screening outcome within the process as an activity
shown in Figure 7.
The process model is then re-evaluated. The included activity (6) is annotated
with a normal effect that realizes Screened(p, ta), and the exceptional effect of
14 George Koliadis et al.
Held(p, ta) is also identified as possibly resulting from the activity. A decision
gateway is applied adjacent to the screening activity to evaluate the outcome
of the activity, and redirect process flow. Upon trajectory decomposition a new
exceptional trajectory is decomposed resulting from the prior addition of the
decision gateway. Goal satisfaction is then re-evaluated. The changes to the
process are successful at resolving the unsatisfied relationship by achieving a
weak satisfaction relationship with two normal and two exceptional trajectories.
Fig. 6. Goal Model Additions for Screening Requirements
Fig. 7. Process Evolution to Achieve Goal Satisfaction
7 Conclusion
We have proposed the GoalBPM methodology that can be used to identify the
satisfaction of a process model against a goal model. The example we presented,
provides a brief and informal overview. There are many possible benefits in applying GoalBPM to current business process design and analysis. This includes
the initial intentional design of business processes that satisfy a deliberate specification of goals. Changes to the business process model may then be made
and tested against the specification of the goals they wish to satisfy in the goal
Supporting BPM with GORE 15
model. Changes may also be made to the goal model and tested against the
current business process model to identify behaviors that are invalidated. Invalid behavior may be explicitly defined, supporting further redesign to align
the processes changed against organizational goals. In order to progress from
the current state, the need for formalism, tool support and testing against a
large, non-trivial business case is required. We are actively pursuing these requirements, which we hope will increase our understanding of the realizability,
workability and viability of GoalBPM for its active use.
References
1. Smith, H., Fingar, P.: Business Process Management: The Third Wave. MeghanKiffer Press, Tampa, FL (2003)
2. Dumas, M., van der Aalst, W.M., ter Hofstede, A.H.: Process-Aware Information Systems: Bridging People and Software Through Process Technology. WileyInterscience (2005)
3. McGoveran, D.: The benefits of a bpms. Technical report, Alternative Technologies, Felton, California, USA (2002)
4. van der Aalst, W., ter Hofstede, A., Weske, M.: Business process management: A
survey. In: BPM’03 - International Conference on Business Process Manag ement,
Berlin, Springer-Verlag, Lecture Notes in Computer Science (2003) 1–12
5. Youngblood, M.D.: Winning cultures for the new economy. Strategy and Leadership 28(6) (2000) 4–9
6. Kavakli, E.: Modelling organizational goals: Analysis of current methods. In:
Proceedings of the 2004 ACM Symposium on Applied Computing, Nicosia, CY
(2004) 1339–1343
7. Pyke, J., Whitehead, R.: Do better maths lead to better business processes? Business Process Trends, http://www.bptrends.com (2004)
8. Wynn, D., Eckert, C., Clarkson, P.J.: Planning business processes in product development organisation s. In: REBPS’03 - Workshop on Requirements Engineering
for Busines s Process Support, Klagenfurt/Veldern, Austria (2003)
9. Watkins, R., Neal, M.: Why and how of requirements tracing. IEEE Software
11(4) (1994) 104–106
10. van Lamsweerde, A.: Goal-oriented requirements engineering: A guided tour. In:
RE’01 - International Joint Conference on Requirements Engi neering, Toronto,
IEEE (2001) 249–263
11. White, S.: Business Process Modeling Notation (BPMN), Version 1.0. Business
Process Management Initiative (BPMI.org). 1.0 edn. (2004)
12. Letier, E., van Lamsweerde, A.: Deriving operational software. In: FSE´ı10 - 10th
ACM SIGSOFT Symp. on the Foundations of S oftware Engineering. (2002)
13. Letier, E.: Reasoning about Agents in Goal-Oriented Requirements Engineerin g.
PhD thesis, Universite Catholique de Louvain, Louvain, Belgium (2001)
14. van Lamsweerde, A., Letier, E.: Handling obstacles in goal-oriented requirements
engineering. IEEE Transactions on Software Engineering 26(10) (2000) 978–1005
15. Bider, I., Johannesson, P.: Tutorial on: Modeling dynamics of business processes
– key for building next generation of business information systems. in: The 21st
international con-ference on conceptual modeling (er2002), tampere, fl, october
7-11, 2002. In: 21st International Conference on Conceptual Modeling (ER2002),
Tampere, FL (2002)
16 George Koliadis et al.
16. Khomyakov, M., Bider, I.: Achieving workflow flexibility through taming the chaos.
In: OOIS’00 - 6th International Conference on Object Oriented I nformation Systems, Springer-Verlag, Berlin (2000) 85–92
17. Andersson, T., Andersson-Ceder, A., Bider, I.: State flow as a way of analysing
business processes - case stud ies. Logistics Information Management 15(1) (2002)
34–45